\section{One specialization with all residual properties}\label{sec:residual}

We now choose a single parameter for the congruence chain.  The geometric
monodromy supplies irreducibility and suitable Frobenius elements, whereas
Hilbert irreducibility allows us to retain these elements after imposing
the local conditions of Section~\ref{sec:local}.  To obtain decomposed
genericity, we must control every place above a split rational prime.
We therefore specialize the two conjugate parameter families together.

Throughout this section, $K,E,p,n,F,U$ have the meanings fixed in
Section~\ref{sec:systems}.  In particular, $p>100$, $n=3p$, and
$F=K(\mu_n)$.  The pairs of covering degree and residual characteristic
that require geometric monodromy are
\begin{equation}\label{res:pairs}
 (m,\ell)=(3,p),\quad(3,5),\quad(p,3).
\end{equation}
In each case $\ell\nmid m$, so the residual cohomology decomposes into
rotation-character spaces after extending its finite coefficient field.

\subsection{Consequences of the two monodromy transformations}

\begin{lemma}\label{res:geometric}
For each pair in \eqref{res:pairs}, the geometric monodromy of
$H^1(C_m,\overline\F_\ell)$ on every nontrivial rotation-character
space is absolutely irreducible.  Its image contains opposite
nonidentity root unipotents on that space.  Moreover, one element
of the geometric monodromy acts regular semisimply on all these
nontrivial character spaces simultaneously.
These statements also hold for the conjugate family over $K$.
\end{lemma}

\begin{proof}
We first work with homology; duality will give the cohomological
statements.  On the character space with generator eigenvalue
$\lambda\in\overline\F_\ell^\times$, the transformations of
Theorem~\ref{mon:transvections} have the form
\[
 T_a=\begin{pmatrix}1&\sigma c_\lambda\\0&1\end{pmatrix},
 \qquad
 T_b=\begin{pmatrix}1&0\\\tau c_\lambda&1\end{pmatrix},
 \qquad
 c_\lambda=2-\lambda-\lambda^{-1},
\]
where $\sigma,\tau\in\{1,-1\}$.  Since
$c_\lambda=-(\lambda-1)^2/\lambda$, it is nonzero when
$\lambda\ne1$.  The two matrices have different unique invariant
lines, even over the algebraic closure.  Hence their common
representation is absolutely irreducible.

The product $T_aT_b^k$ has determinant one and trace
\begin{equation}\label{res:trace-product}
 2+\sigma\tau k c_\lambda^2.
\end{equation}
In odd characteristic, a determinant-one matrix is regular
semisimple exactly when its trace is different from $2$ and $-2$.
If $m=3$, then $\lambda+\lambda^{-1}=-1$, so $c_\lambda=3$ for both
nontrivial characters.  Take $k=2$.  The two possible traces in
\eqref{res:trace-product} are $20$ and $-16$.  Modulo $5$ these
are $0$ and $4$, and neither equals $2$ or $-2$; modulo $p>100$
they again avoid $2$ and $-2$.

For $m=p$ and $\ell=3$, take $k=1$.  Failure of regular
semisimplicity would force $c_\lambda^2=4$ or $-4$, since
$c_\lambda\ne0$.  Multiplying either equation by $\lambda^2$
gives a polynomial of degree four over $\F_3$ satisfied by
$\lambda$.  Its degree over $\F_3$ is therefore some $d\le4$,
and its order divides $3^d-1\le80$.  This contradicts the fact
that its order is the prime $p>100$.  Thus the same product works
for every nontrivial character in this case as well.

The transformations come from loops in the constant-elliptic-curve
parameter space.  As explained in Theorem~\ref{mon:transvections},
that family is a pullback of the family over $U$.  Its monodromy
image is therefore contained in the geometric monodromy image
over $U$.  Duality preserves absolute irreducibility, nonidentity
unipotents, and regular semisimplicity.  This proves the assertions
for cohomology.  The geometric argument applies to every complex
elliptic curve, so it applies equally to the conjugate of $E$.
\end{proof}

\subsection{The original elliptic curve and its conjugate}

The trivial character requires a different argument, since the
geometric family varies only a quadratic twist of $E$.  Its large
residual image was built into the choice of $p$.  The following
count shows that no independence assumption between $E$ and its
conjugate is necessary.

\begin{lemma}\label{res:trivial-character}
For every $t\in U(K)$, the representation
$\overline V_{1,p}$ is absolutely irreducible on $G_{F(\zeta_p)}$
and is decomposed generic over $F$.
\end{lemma}

\begin{proof}
Let $\Gamma$ be the image of $G_K$ on $E[p]$.
The abelian extension $F/K$ implies that the image of $G_F$
contains $[\Gamma,\Gamma]$.  Since $\Gamma$ contains
$\SL_2(\F_p)$ and this group is perfect for $p>3$, its commutator
subgroup is $\SL_2(\F_p)$: one inclusion follows from
perfectness, and the other from the determinant map.
Moreover, the determinant is trivial on $G_F$ because
$\mu_p\subset F$.  We have therefore proved
\begin{equation}\label{res:SL2-image}
 \rho_{E,p}(G_F)=\SL_2(\F_p).
\end{equation}
The same equality holds for the conjugate curve $E^c$, where $c$
is the nontrivial automorphism of $K/\Q$, because $F/\Q$ is normal.

Consider the joint image
\[
 J\subset\SL_2(\F_p)\times\SL_2(\F_p)
\]
of $G_F$ on $E[p]$ and $E^c[p]$.  Both projections are surjective.
In $\SL_2(\F_p)$ there are exactly $p^2$ matrices of trace $2$.
Indeed, writing a matrix as $\left(\begin{smallmatrix}a&b\\c&d
\end{smallmatrix}\right)$, the trace and determinant conditions
give $d=2-a$ and $bc=-(a-1)^2$.  The value $a=1$ gives $2p-1$
choices for $(b,c)$; the other $p-1$ values of $a$ each give
$p-1$ choices.  Their sum is $p^2$.  Replacing a matrix by its
negative gives the same count for trace $-2$.

The order of $\SL_2(\F_p)$ is $p(p^2-1)$.  Since each projection
of $J$ has equal-sized fibers, the proportion of elements of $J$
whose image in at least one factor has trace $2$ or $-2$ is at most
\[
 2\frac{2p^2}{p(p^2-1)}
 =\frac{4p}{p^2-1}<1.
\]
Thus some element of $G_F$ acts regular semisimply on both
$E[p]$ and $E^c[p]$.

Take a finite Galois extension of $\Q$ containing $F$ and the
fields needed for these two residual actions.  Conjugation by
$G_\Q$ either preserves the two curves or exchanges them, and
conjugates their residual matrices.  The selected element and all
its conjugates therefore act regular semisimply on the
representation of $E$.  Chebotarev supplies rational primes $q$
splitting completely in $F$ with this property at every $v\mid q$
of $F$.  We exclude $p$ and the finitely many primes of bad
reduction.  Since $\mu_p\subset F$, these primes satisfy
$q\equiv1\pmod p$.  Distinct Frobenius eigenvalues then give the
ratio condition in the definition of decomposed genericity.

The curve $B_t$ is a quadratic twist of $E$.  Twisting multiplies
residual matrices by scalar signs and hence preserves absolute
irreducibility and eigenvalue ratios.  Excluding the finitely many
ramified primes of this twist leaves the same Chebotarev argument
available.  Finally, $\zeta_p\in F$, and taking the dual of $E[p]$
gives the residual $H^1$ realization.  Proposition~\ref{sys:elliptic}
now proves both assertions for $\overline V_{1,p}$.
\end{proof}

\subsection{Simultaneous specialization}

We have obtained the requisite geometric elements and dealt with
the trivial character.  It remains to retain the geometric elements
at one parameter and arrange that their Frobenius properties hold
at all places above a rational prime.  Restriction of scalars makes
the two conjugate parameter families visible in the same finite
Galois cover.

We next record precisely the specialization theorem needed to
preserve these finite monodromy images.  Allowing a constant field
in the cover is essential: the Frobenius primes used below must
split in a prescribed field.

\begin{lemma}\label{res:finite-specialization}
Let $W$ be a nonempty open subset of $\A^d_\Q$, let $D/\Q$ be a
finite Galois extension, and let $H$ be a finite quotient of
$\pi_1(W)$ through which the constant-field map factors as
\[
 \pi_1(W)\twoheadrightarrow H\twoheadrightarrow\Gal(D/\Q).
\]
Prescribe nonempty open subsets of
$W(\Q_r)$ for finitely many rational primes $r$.
There is a point $x\in W(\Q)$ in all the prescribed opens for
which the specialization map $s_x:G_\Q\to H$ is surjective.
For this point,
\begin{equation}\label{res:constant-kernel}
 s_x(G_D)=\ker\bigl(H\longrightarrow\Gal(D/\Q)\bigr).
\end{equation}
\end{lemma}

\begin{proof}
The quotient $H$ defines a connected finite \(\acute etale\) Galois
cover of $W$.  Hilbert irreducibility with weak approximation gives
a rational point in the prescribed local opens for which its fiber
is connected; see \cite[Proposition~3.3.1 and
Theorem~3.5.3]{SerreHilbert}.  Connectedness of this Galois fiber is
equivalent to surjectivity of $s_x$.  Geometric connectedness of
the cover is not required.  The composition of $s_x$ with
$H\to\Gal(D/\Q)$ is the usual restriction map of absolute Galois
groups.  Consequently, an element of $H$ in the displayed kernel
has a preimage in $G_D$, which proves
\eqref{res:constant-kernel}.
\end{proof}

\begin{proposition}\label{res:specialization}
Let $S_K$ be a finite set of finite places of $K$ containing all
places above $3p$.  For each $v_0\in S_K$, prescribe a nonempty
open subset $\Omega_{v_0}\subset U(K_{v_0})$.  There exists a
single $t\in U(K)$ lying in every $\Omega_{v_0}$ such that the
following assertions hold over $F=K(\mu_{3p})$.
\begin{enumerate}
\item For each $(m,\ell)$ in \eqref{res:pairs} and each character
$\theta$ of order $m$, the representation
$\overline V_{\theta,\ell}$ is absolutely irreducible on
$G_{F(\zeta_\ell)}$ and is decomposed generic.
\item The representation $\overline V_{1,p}$ has these same two
properties at $p$.
\item The elliptic curve $A_t=\Jac(C_3(t)/\langle w\rangle)$ is
non-CM.
\end{enumerate}
In particular, when the opens above $3p$ are those of
Proposition~\ref{loc:opens}, this same parameter also satisfies
all the local conclusions of Corollary~\ref{loc:characters}.
\end{proposition}

\begin{proof}
Put
\[
 W=\Res_{K/\Q}U.
\]
Since $U$ is an open subset of $\A^1_K$, the scheme $W$ is an open
subset of $\A^2_\Q$.  Choose a finite Galois extension $D/\Q$
containing $F$ and $\zeta_\ell$ for every residual characteristic
in \eqref{res:pairs}.  Let $c$ be the nontrivial automorphism of
$K/\Q$, and write $U^c$ and $\mathcal C_m^c$ for the conjugate
parameter open and curve family.  Then
\[
 W_D\simeq U_D\times (U^c)_D.
\]
For all three pairs, pull back the finite local system
$H^1(C_m,\F_\ell)$ from the first factor and
$H^1(C_m^c,\F_\ell)$ from the second.

Their joint monodromy has finite image on $\pi_1(W_D)$.  Take the
normal core, in $\pi_1(W)$, of the common kernel of these actions.
It is an open normal subgroup contained in $\pi_1(W_D)$.
The resulting finite quotient $H$ of $\pi_1(W)$ records every
listed action and has a quotient $\Gal(D/\Q)$.

Over $\C$, a loop in either parameter factor acts trivially on
local systems pulled back from the other factor.  Thus the
geometric monodromy transformations of Lemma~\ref{res:geometric}
are available independently in both factors.  In particular, for
each fixed pair $(m,\ell)$ we can choose an element $h_{m,\ell}$
of the image of geometric monodromy that is regular semisimple
on every nontrivial character space in both factors.  To obtain
it, take a product of the two factor loops furnished by that
lemma.  Every such element maps trivially to $\Gal(D/\Q)$.
The quotient also contains the individual opposite unipotents
on every character space.

The prescribed local conditions are compatible with this
restriction of scalars.  For a rational prime $r$ one has
\[
 W(\Q_r)=\prod_{v_0\mid r}U(K_{v_0}).
\]
The sets $\Omega_{v_0}$ therefore specify nonempty open subsets
of $W(\Q_r)$ for the finitely many relevant $r$; at unprescribed
factors impose no additional condition.  Apply
Lemma~\ref{res:finite-specialization}.  It gives a point
$x\in W(\Q)$, equivalently $t\in U(K)$, satisfying all these
local conditions and inducing a surjection $G_\Q\to H$.
Over $D$, this point has coordinates $(t,t^c)$ in the displayed
product. Smooth proper base change identifies the specialized
local systems with the residual cohomology of $C_m(t)$ and
$C_m^c(t^c)$.

By \eqref{res:constant-kernel}, every geometric monodromy element
used above is the image of an element of $G_D$.  The opposite
unipotents therefore survive on each specialized character space
already over $G_D$, and prove its absolute irreducibility.
Since $G_D\subset G_{F(\zeta_\ell)}$, this also proves absolute
irreducibility on the latter group.  Because $\ell\nmid m$,
taking rotation eigenspaces in integral cohomology of $C_m$
commutes with reduction after adjoining the character values.
These specialized character spaces are thus reductions of
$V_{\theta,\ell}$.  The rational identifications and independence
of semisimplified reduction in Section~\ref{sec:systems} show
that they give precisely $\overline V_{\theta,\ell}$, including
when this system is defined using $C_n$.

To establish decomposed genericity, fix a pair $(m,\ell)$ and
choose a preimage $g\in G_D$ of $h_{m,\ell}$.
An element $\sigma\in G_\Q$ preserves or
exchanges the two specialized covers $C_m(t)$ and $C_m^c(t^c)$,
according to its restriction to $K$; it also permutes the rotation
characters through its action on the roots of unity.  The action
of $\sigma g\sigma^{-1}$ on a character space of the first cover
is therefore conjugate to the action of $g$ on some character
space of one of the two covers.
The construction imposed regular semisimplicity on all these
spaces, so the property holds for every such conjugate.

The kernel of the surjective specialization $G_\Q\to H$ defines
a finite Galois extension of $\Q$.  It contains $D$.
Chebotarev applied to the conjugacy class of $h_{m,\ell}$ in
this extension gives rational primes $q$ whose Frobenius class
is that class.  They split completely in $D$, since
$h_{m,\ell}$ is trivial on $D$.  Discard $q=\ell$ and the finitely
many primes where the curves or residual representations are
ramified.  For every $v\mid q$ of $F$, its Frobenius is represented
by a conjugate of the selected element, and acts regular
semisimply on each required character space.  Here the residue
degree is one because $q$ splits in $F$.  Moreover,
$q\equiv1\pmod\ell$ because $\zeta_\ell\in D$.
Thus the ratio of the two Frobenius eigenvalues is different from
$q$ in characteristic $\ell$, at every $v\mid q$.  This is
exactly decomposed genericity.  A different $q$ may be used for
each pair, as the definition permits.  Assertion (2) follows
from Lemma~\ref{res:trivial-character}.

It remains to exclude CM for $A_t$.  For $(m,\ell)=(3,p)$ the
preceding specialization places two opposite nonidentity
unipotents in the residual image of $V_{\eta,p}$.
Proposition~\ref{sys:elliptic} identifies this with the residual
cohomology of $A_t$: both are irreducible, so their common
semisimplification is the representation itself.
If $A_t$ had CM, all its endomorphisms would be defined over an
extension of $F$ of degree at most two.  Over that extension,
the characteristic-zero Galois image lies in the commutative
centralizer of the CM algebra.  Its reduction is abelian as well.
Consequently the full residual image would have an abelian
normal subgroup of index at most two.

Each of the two unipotents has the odd prime order $p$, so both
would belong to that abelian subgroup.  But
\[
 \begin{pmatrix}1&a\\0&1\end{pmatrix}
 \begin{pmatrix}1&0\\b&1\end{pmatrix}
 \ne
 \begin{pmatrix}1&0\\b&1\end{pmatrix}
 \begin{pmatrix}1&a\\0&1\end{pmatrix}
 \qquad(a b\ne0).
\]
This contradiction proves (3).  Finally, the local assertions
follow by choosing the prescribed opens as in
Proposition~\ref{loc:opens}; the Hilbert specialization was made
inside all those opens.
\end{proof}
